% META: {"id":"1NSI-TC-EX-032","chapitre":"1NSI-TYPES-CONSTRUITS","type_objet":"exercice","capacites":["1NSI-TYPES-CONSTRUITS-C3"],"parcours":2,"competences":["analyser"],"duree_min":10,"mode_creation":"ex_nihilo","sources_inspiration":[],"status":"approved","origin":"needs_review","status_history":["needs_review","audited","accepted","approved"],"release_acceptance":"RELEASE_OWNER_BATCH_ACCEPTANCE","acceptance_closure_digest":"sha256:9b3ccf9a81c5520fbb7e03b7b2d3e7bf2057a02834b6d908bf3be5f49f3b553f","acceptance_receipt_digest":"sha256:4fad86f0282e0fa4e0419cca1ad6379ae7af5a280c04a4a90880f90a1c044242"}
% BEGIN-TRACE
% scores = [[10, 8, 12],
%            [15, 11, 9],
%            [7, 14, 13]]
% moyennes = []
% for ligne in scores:
%     s = 0
%     for v in ligne:
%         s = s + v
%     moyennes.append(s / len(ligne))
% print(moyennes)
% EXPECTED
% [10.0, 11.666666666666666, 11.333333333333334]
% END-TRACE
\begin{exercice}{1NSI-TC-EX-032}{2}{10}
Trois équipes d'e-sport ont chacune disputé trois matchs. Leurs scores sont
enregistrés dans la grille suivante :
\begin{python}
scores = [[10, 8, 12],
           [15, 11, 9],
           [7, 14, 13]]
moyennes = []
for ligne in scores:
    s = 0
    for v in ligne:
        s = s + v
    moyennes.append(s / len(ligne))
print(moyennes)
\end{python}
\begin{enumerate}
  \item Donner, sans exécuter, l'affichage produit.
  \item Modifier le code pour n'afficher que la moyenne la plus élevée.
\end{enumerate}
\end{exercice}
